\chapter{Мембранный потенциал}
\label{sec:membrane}

\section{Откуда берётся напряжение}

Воткнём тонкий стеклянный электрод в нейрон, а второй электрод положим в раствор рядом с клеткой. Вольтметр покажет около $-70\mV$: внутри клетки потенциал ниже, чем снаружи. Это напряжение называется \emph{потенциалом покоя}. Откуда оно берётся? У клетки нет ни генератора, ни химической батареи в обычном смысле. Есть только два раствора с разным составом, разделённые плёнкой, в которой есть поры. Оказывается, этого достаточно.

Начнём с мысленного опыта. Пусть по обе стороны мембраны --- раствор хлорида калия, внутри концентрированный, $140\mM$, снаружи разбавленный, $5\mM$. Пусть сначала мембрана вообще непроницаема, и оба раствора электронейтральны: потенциалы внутри и снаружи равны. Теперь откроем в мембране поры, пропускающие \emph{только} калий. Ионы калия начнут диффундировать туда, где их меньше, --- наружу. Каждый вышедший ион уносит положительный заряд, а оставшийся внутри отрицательный ион хлора выйти не может. Внутренняя сторона мембраны заряжается отрицательно, внешняя --- положительно. Возникшее поле тянет ионы калия обратно внутрь. Чем больше ионов вышло, тем сильнее поле, и в какой-то момент электрическая сила уравновешивает диффузионную: поток наружу равен потоку внутрь. Напряжение, при котором это происходит, называется \emph{равновесным потенциалом} для калия.

Важно, сколько ионов для этого нужно. В главе~\ref{sec:neuron} мы оценили: чтобы зарядить мембрану до десятков милливольт, достаточно сдвинуть примерно стотысячную долю ионов. Концентрации при этом практически не меняются. Поэтому мы можем считать концентрации внутри и снаружи заданными, а потенциал --- определяемым ими.

\section{Формула Нернста}

Найдём равновесный потенциал количественно. Сделаем это двумя способами: через потоки и через больцмановское распределение.

\emph{Через потоки.} Поток ионов сорта $i$ через единицу площади складывается из диффузионного и дрейфового:
\begin{equation}
J_i \;=\; -D_i\,\frac{\dd c_i}{\dd x} \;-\; D_i\,\frac{z_iF}{RT}\,c_i\,\frac{\dd\varphi}{\dd x} .
\label{eq:nernstplanck}
\end{equation}
Первое слагаемое --- закон Фика: ионы текут туда, где их меньше. Второе --- дрейф в электрическом поле $-\dd\varphi/\dd x$: подвижность иона связана с коэффициентом диффузии соотношением Эйнштейна, отсюда множитель $D_i z_iF/RT$. Здесь $z_i$ --- заряд иона в единицах элементарного ($+1$ для калия, $+2$ для кальция, $-1$ для хлора), $F\approx 96\,485$~Кл/моль --- постоянная Фарадея, $R$ --- газовая постоянная, $T$ --- абсолютная температура. Уравнение~\eqref{eq:nernstplanck} называется уравнением Нернста---Планка.

В равновесии $J_i=0$. Разделив на $D_ic_i$, получаем
\[
\frac{\dd\ln c_i}{\dd x} \;=\; -\frac{z_iF}{RT}\,\frac{\dd\varphi}{\dd x},
\]
и, проинтегрировав поперёк мембраны от внешней стороны до внутренней,
\[
\ln\frac{c_i^{\text{вн}}}{c_i^{\text{нар}}} \;=\; -\frac{z_iF}{RT}\,\bigl(\varphi_{\text{вн}}-\varphi_{\text{нар}}\bigr).
\]
Разность потенциалов в скобках --- это и есть мембранный потенциал. Обозначив равновесное значение через $E_i$, получаем формулу Нернста:
\begin{empheq}[box=\keybox]{equation}
E_i \;=\; \frac{RT}{z_iF}\,\ln\frac{\conc{i}_{\text{нар}}}{\conc{i}_{\text{вн}}} .
\label{eq:nernst}
\end{empheq}
Обратите внимание: ни коэффициент диффузии, ни толщина мембраны, ни то, как устроены поры, в ответ не вошли. Равновесие --- свойство термодинамическое, а не кинетическое.

\emph{Через Больцмана.} Тот же ответ можно получить за одну строчку. В равновесии вероятность найти ион в точке с потенциалом $\varphi$ пропорциональна $\exp(-z_ie\varphi/kT)$, поэтому отношение концентраций по обе стороны мембраны равно $\exp(-z_ie\,V/kT)$. Прологарифмировав и заметив, что $e/k=F/R$, получаем~\eqref{eq:nernst}. Этот вывод полезен тем, что показывает смысл величины $kT/e$: это напряжение, при котором электрическая энергия одного элементарного заряда равна тепловой энергии.

Посчитаем её. При температуре тела $T=310$~К
\begin{empheq}[box=\keybox]{equation}
\frac{RT}{F} \;=\; \frac{kT}{e} \;\approx\; 26.7\mV .
\label{eq:thermalvoltage}
\end{empheq}
Для калия с концентрациями из таблицы~\ref{tab:concentrations} получаем $E_{\rm K}=26.7\cdot\ln(5/140)\approx-89\mV$. Так же для остальных ионов --- см. таблицу~\ref{tab:nernst}.

\begin{table}[t]
\centering
\begin{tabular}{lccc}
\toprule
Ион & $z$ & $[\,\cdot\,]_{\text{нар}}/[\,\cdot\,]_{\text{вн}}$ & $E_i$, мВ \\
\midrule
$\K$  & $+1$ & $5/140$ & $-89$ \\
$\Na$ & $+1$ & $145/12$ & $+67$ \\
$\Cl$ & $-1$ & $110/7$ & $-74$ \\
$\Ca$ & $+2$ & $2/10^{-4}$ & $+132$ \\
\bottomrule
\end{tabular}
\caption{Равновесные потенциалы, посчитанные по формуле Нернста~\eqref{eq:nernst} при $T=310$~К для концентраций таблицы~\ref{tab:concentrations}. Полезное правило: изменение отношения концентраций в десять раз сдвигает $E$ на $61\mV$ для одновалентного иона.}
\label{tab:nernst}
\end{table}

\section{Несколько ионов: формула Гольдмана---Ходжкина---Каца}

Потенциал покоя, $-70\mV$, не совпадает ни с одним из равновесных потенциалов таблицы~\ref{tab:nernst}. Он близок к калиевому, но немного выше. Причина в том, что мембрана в покое пропускает не только калий. Она в основном калиевая, но через неё немного подтекает натрий, который тянет потенциал к $+67\mV$, и хлор. В итоге устанавливается не равновесие, а стационарное состояние, в котором \emph{суммарный} ток через мембрану равен нулю, хотя каждый ион по отдельности течёт.

Чтобы найти этот потенциал, надо решить уравнение Нернста---Планка~\eqref{eq:nernstplanck} при $J_i\ne0$. Для этого нужна дополнительная гипотеза о профиле поля внутри мембраны. Простейшая из них --- поле внутри мембраны постоянно~\cite{Goldman1943}. Тогда уравнение для каждого иона линейно и решается явно. Приравняв нулю сумму токов калия, натрия и хлора, получаем формулу Гольдмана---Ходжкина---Каца~\cite{HodgkinKatz1949}:
\begin{empheq}[box=\keybox]{equation}
V_{\text{покоя}} \;=\; \frac{RT}{F}\,\ln\frac{P_{\rm K}\conc{K^+}_{\text{нар}}+P_{\rm Na}\conc{Na^+}_{\text{нар}}+P_{\rm Cl}\conc{Cl^-}_{\text{вн}}}{P_{\rm K}\conc{K^+}_{\text{вн}}+P_{\rm Na}\conc{Na^+}_{\text{вн}}+P_{\rm Cl}\conc{Cl^-}_{\text{нар}}} .
\label{eq:ghk}
\end{empheq}
Здесь $P_i$ --- проницаемости мембраны для каждого иона. Обратите внимание, что у хлора, отрицательного иона, внешняя и внутренняя концентрации поменялись местами. Формула разумно ведёт себя в предельных случаях: если проницаема только одна разновидность ионов, она переходит в формулу Нернста для неё.

Подставим типичные для покоя отношения проницаемостей $P_{\rm K}:P_{\rm Na}:P_{\rm Cl}=1:0.05:0.45$. Числитель: $5+0.05\cdot145+0.45\cdot7\approx15.4$. Знаменатель: $140+0.05\cdot12+0.45\cdot110\approx190$. Получаем $V=26.7\cdot\ln(15.4/190)\approx-67\mV$, в хорошем согласии с опытом.

\begin{keyblock}
\begin{proposition}[Мембранный потенциал как голосование]
\label{prop:vote}
Мембранный потенциал всегда лежит между наименьшим и наибольшим равновесным потенциалом тех ионов, которые проходят через мембрану. Каждый ион тянет потенциал к своему $E_i$ с силой, пропорциональной проницаемости мембраны для него. Изменить мембранный потенциал можно двумя способами: изменить концентрации, то есть сами $E_i$, или изменить проницаемости, то есть веса голосов.
\end{proposition}
\end{keyblock}

Нейрон почти всегда пользуется вторым способом. Концентрации держат насосы, и они меняются медленно. А проницаемости меняются за доли миллисекунды: открылись натриевые каналы --- и голос натрия стал громче всех, потенциал рванулся к $+67\mV$. Это и есть потенциал действия, которому посвящена следующая глава.

\section{Эквивалентная схема мембраны}

Для всего дальнейшего удобно заменить мембрану электрической схемой. Каждый сорт ионов, проходящих через мембрану, --- это ветвь из батарейки с ЭДС $E_i$ и последовательного с ней сопротивления, или, что удобнее, проводимости $g_i$. Ток через ветвь по закону Ома равен
\begin{empheq}[box=\keybox]{equation}
I_i \;=\; g_i\,(V-E_i) .
\label{eq:ohm}
\end{empheq}
Разность $V-E_i$ называется \emph{движущей силой} для иона $i$. Знаки выбраны так, что положительный ток --- это выход положительного заряда из клетки. Например, при $V=-70\mV$ движущая сила для натрия равна $-137\mV$, ток отрицательный: натрий входит в клетку. Параллельно всем ветвям стоит ёмкость мембраны $C$ (рис.~\ref{fig:circuit}).

\begin{figure}[t]
\centering
\begin{circuitikz}[scale=0.95]
\draw (0,4) -- (9,4);
\draw (0,0) -- (9,0);
\draw (0,4) to[C,l=$C_m$] (0,0);
\draw (2.5,4) to[R,l=$g_{\rm K}$] (2.5,2) to[battery1,l=$E_{\rm K}$] (2.5,0);
\draw (5,4) to[vR,l=$g_{\rm Na}$] (5,2) to[battery1,l=$E_{\rm Na}$] (5,0);
\draw (7.5,4) to[R,l=$g_{L}$] (7.5,2) to[battery1,l=$E_{L}$] (7.5,0);
\draw (9,4) to[I,l=$I_{\text{ext}}$] (9,0);
\node[above] at (4.5,4.05) {\small внутри клетки};
\node[below] at (4.5,-0.05) {\small снаружи};
\end{circuitikz}
\caption{Эквивалентная схема участка мембраны. Ёмкость $C_m$ --- липидный слой; каждая ветвь с батарейкой --- ионные каналы одного сорта. Стрелка на натриевой проводимости напоминает, что она зависит от напряжения (глава~\ref{sec:actionpotential}); проводимость утечки $g_L$ собирает всё остальное. Внешний ток $I_{\text{ext}}$ --- то, что вводит экспериментатор через электрод, или синаптический ток.}
\label{fig:circuit}
\end{figure}

Закон Кирхгофа для узла «внутри клетки» говорит, что ток, пришедший извне, идёт либо на зарядку ёмкости, либо через ионные ветви:
\begin{empheq}[box=\keybox]{equation}
C\,\frac{\dd V}{\dd t} \;=\; -\sum_i g_i\,(V-E_i) \;+\; I_{\text{ext}} .
\label{eq:membrane}
\end{empheq}
Это уравнение --- основа всей количественной нейрофизиологии. В главе~\ref{sec:actionpotential} мы добавим в него зависимость $g_i$ от напряжения, в главе~\ref{sec:synapse} --- синаптические проводимости, в главе~\ref{sec:cable} --- пространство. Но уже сейчас из него можно извлечь два важных вывода.

\emph{Первый: потенциал покоя.} При $I_{\text{ext}}=0$ и постоянных проводимостях стационарное решение~\eqref{eq:membrane} --- взвешенное среднее:
\begin{equation}
V_{\text{покоя}} \;=\; \frac{\sum_i g_iE_i}{\sum_i g_i} .
\label{eq:weighted}
\end{equation}
Это то же «голосование», что и в утверждении~\ref{prop:vote}, только веса теперь --- проводимости, а не проницаемости. Формула~\eqref{eq:weighted} проще, чем~\eqref{eq:ghk}, и для большинства целей её достаточно.

\emph{Второй: постоянная времени.} Обозначим $g=\sum_i g_i$ полную проводимость. Если ток $I_{\text{ext}}$ включить скачком, напряжение приближается к новому стационарному значению по экспоненте
\begin{equation}
V(t) \;=\; V_\infty + (V_0-V_\infty)\,e^{-t/\tau}, \qquad \tau = \frac{C}{g} = \frac{c_m}{g_m} ,
\label{eq:tau}
\end{equation}
где $c_m$ и $g_m$ --- ёмкость и проводимость единицы площади мембраны. Площадь сокращается, и постоянная времени оказывается свойством мембраны, а не клетки. Удельное сопротивление мембраны покоящегося нейрона $r_m=1/g_m$ составляет около $10$--$50$~кОм$\cdot$см$^2$. Отсюда $\tau=c_mr_m\approx10^{-6}\cdot2\cdot10^4=20\ms$. Это время, за которое нейрон «забывает» полученный толчок. Если два входных сигнала приходят с интервалом меньше $\tau$, они успевают сложиться; если больше --- первый уже рассосался, когда пришёл второй. Постоянная времени задаёт окно, в котором нейрон интегрирует свои входы.

\section{Насос как зарядное устройство}

Мы говорили, что концентрации держат насосы. На языке схемы ионные батарейки $E_i$ --- аккумуляторы, а насос --- зарядное устройство, которое непрерывно их подзаряжает. Но в стационарном состоянии через мембрану непрерывно текут натрий внутрь и калий наружу. Если бы насосы остановились, перепады концентраций постепенно выровнялись бы, и потенциал покоя исчез бы. Сколько на это нужно времени? Утечки малы, а запас ионов велик, поэтому в покое клетка без насоса разрядилась бы за минуты или часы. Но активный нейрон, выдающий десятки импульсов в секунду, расходует запас намного быстрее, и насосы должны работать непрерывно. Мы уже видели в главе~\ref{sec:neuron}, что на это уходит большая часть энергии мозга. Если питание зарядного устройства отключить, батарейки медленно садятся, и через несколько минут схема перестаёт работать.

Натрий-калиевый насос выносит три положительных заряда, а вносит два. Значит, он сам создаёт небольшой ток наружу и немного понижает потенциал, на несколько милливольт. Такой насос называют \emph{электрогенным}. Но главная его роль --- не этот прямой вклад, а поддержание перепадов концентраций, без которых не было бы ни потенциала покоя, ни импульса.
